\documentclass[11pt,a4paper]{article}
\usepackage[margin=1in]{geometry}
\usepackage[T1]{fontenc}
\usepackage{lmodern}
\usepackage{microtype}
\usepackage{hyperref}
\usepackage{array}
\usepackage{longtable}
\hypersetup{
  colorlinks=true,
  linkcolor=black,
  urlcolor=black,
  pdftitle={OC Core 1.3.2 Expert Technical Spine},
  pdfauthor={Alexander Yashin},
  pdfsubject={Ontology of Continua Core 1.3.2 expert technical spine},
  pdfkeywords={Ontology of Continua, OC Core 1.3.2, technical spine}
}
\title{\textbf{OC Core 1.3.2 Expert Technical Spine}\\
\large Formal Navigation for Reviewers of the Release Candidate Package}
\author{Alexander Yashin\\Independent Researcher, Leipzig/Halle, Germany\\ORCID 0009-0008-6166-0914}
\date{Version v1.3.2; release date: 28.04.2026}
\begin{document}
\maketitle

\section*{Technical Role}
This spine gives a compact expert route through OC Core 1.3.2. It is meant for reviewers who want to inspect the formal and evidential architecture before reading the entire master monograph.

\section*{Core Grammar}
OC Core represents a continuum through identity-bearing conditions, admissible realizations, thresholds, embeddings, and history-sensitive classification. The release-relevant grammar distinguishes:
\begin{itemize}
  \item live continuation under preserved admissible realization;
  \item death or collapse when admissible realization is empty;
  \item residue after death without live identity;
  \item rebirth when later live structure is grounded in residue without being identical to the original continuum.
\end{itemize}

\section*{Claim-Support Classes}
\begin{longtable}{>{\raggedright\arraybackslash}p{0.34\linewidth}>{\raggedright\arraybackslash}p{0.56\linewidth}}
\textbf{Class} & \textbf{Release meaning}\\
\hline
FORMALLY\_PROVED & The release contains or references a proof route adequate for the claim scope.\\
THEOREM\_NATIVE\_HELD\_OUT\_VALIDATED & The theorem label is retained only where bounded proof support exists.\\
OPERATIONALLY\_SUPPORTED\_WITHIN\_BOUNDS & Empirical or domain support is present, but the claim remains bounded by protocol and falsifier.\\
SIMULATION\_ILLUSTRATION\_ONLY & Simulation clarifies behavior without serving as standalone proof.\\
PUBLIC\_ROUTE\_DISCOVERY\_ONLY & Source route is exposed for audit, not used as strong support.\\
FRONTIER\_WORK & Future research surface; not a canonical Core promotion.\\
\end{longtable}

\section*{Release Technical Invariants}
\begin{itemize}
  \item No canonical claim is promoted by a support packet alone.
  \item No prediction claim is accepted without replay protocol, code, fixture or data, output hash, tolerance envelope, and falsifier.
  \item No external criticism response is treated as closure unless release-visible evidence, proof route, replay, or owner-gate terminal state exists.
  \item No private reviewer-source material or private model transcript is included in the public package.
  \item No publication action is lawful while the no-send lock remains true.
\end{itemize}

\section*{Primary Artifact Map}
\begin{longtable}{>{\raggedright\arraybackslash}p{0.32\linewidth}>{\raggedright\arraybackslash}p{0.58\linewidth}}
\textbf{Artifact} & \textbf{Expert use}\\
\hline
Master monograph & Full formal narrative, definitions, proof route, domain appendices, and reviewer navigation.\\
Journal core & Compact source-audited theory kernel.\\
Methods companion & Build, replay, source-audit, and release-governance route.\\
Critique map & Objection families, response surfaces, and blocker ledger reading order.\\
Release dossier & Artifact inventory, gates, no-send lock, and owner action route.\\
Science blocker closure ledger & Row-level terminality of the release blockers.\\
\end{longtable}

\section*{Owner Review Boundary}
OC Core 1.3.2 can be reviewed as a coherent release package. It still cannot publish itself. Owner approval, publication channel decision, and final artifact freeze confirmation are separate lawful actions.

\section*{Reviewer Routes}
\begin{longtable}{>{\raggedright\arraybackslash}p{0.25\linewidth}>{\raggedright\arraybackslash}p{0.65\linewidth}}
\textbf{Reviewer} & \textbf{Technical route}\\
\hline
Owner & No-send lock, owner packet, package freeze hash, release dossier, and publication risk memo.\\
Mathematical reviewer & Formal definitions, theorem roadmap, proof machinery appendix, closure ledger, and support map.\\
Domain reviewer & K/domain projection matrix, replay manifests, DRT strict salvage, and domain packet audit.\\
Reproducibility reviewer & Manifest, checksums, ZIP integrity, data/simulation manifests, and PDF text-quality scan.\\
Publication reviewer & LRGEF, Parfitian Cerberus, citation metadata, release notes, and postflight checklist.\\
\end{longtable}

\section*{Proven, Replayed, Future}
\begin{longtable}{>{\raggedright\arraybackslash}p{0.28\linewidth}>{\raggedright\arraybackslash}p{0.58\linewidth}}
\textbf{Class} & \textbf{Technical interpretation}\\
\hline
Proven or evidence-closed & Release support is anchored in definitions, theorem-fate rows, claim-ledger support, or the terminal closure ledger.\\
Replayed & Release support is anchored in deterministic package artifacts, data or fixture manifests, output hashes, and tolerance/falsifier records.\\
Future work & The item is deliberately outside 1.3.2 authority and may only become stronger through a later release gate.\\
\end{longtable}

\section*{Support-Map Discipline}
The expert reader should use the prediction and promotion support map as a guardrail. It does not make weak statements strong. It binds risky words to their allowed support route: proof, replay, claim ledger, structural-only boundary, or future-work route. Any statement that cannot be bound must be rewritten before release.

\section*{Technical Acceptance Criteria}
\begin{itemize}
  \item six primary PDFs have current identity, adequate page count, and clean extracted text;
  \item all eighty-two closure rows carry row-level support class, source anchor, evidence refs, manuscript refs, release refs, and terminal rationale;
  \item the ZIP contains no unclassified legacy primary PDFs or nonpublic template payloads;
  \item every strong prediction or promoted-claim occurrence is represented in the support map;
  \item public release remains impossible until owner/channel actions are explicitly unlocked.
\end{itemize}

\clearpage
\section*{Technical Appendix A: Formal Review Table}
\begin{longtable}{>{\raggedright\arraybackslash}p{0.24\linewidth}>{\raggedright\arraybackslash}p{0.66\linewidth}}
\textbf{Formal surface} & \textbf{Question to ask}\\
\hline
Definitions & Are identity, admissible realization, collapse, residue, and rebirth defined before use?\\
Theorem route & Is every theorem-native phrase bound to a theorem-fate or proof-support row?\\
K-level system & Does each K-level projection state match the domain and inclusion ledgers?\\
Claim ledger & Does each promoted claim carry an allowed support class and evidence reference?\\
Release machine & Does the gate suite pass without enabling publication?\\
\end{longtable}

\clearpage
\section*{Technical Appendix B: Evidence Review Table}
\begin{longtable}{>{\raggedright\arraybackslash}p{0.24\linewidth}>{\raggedright\arraybackslash}p{0.66\linewidth}}
\textbf{Evidence type} & \textbf{Sufficiency standard}\\
\hline
Proof route & Source-audited formal statement, theorem-fate classification, and clear release scope.\\
Replay & Fixed inputs, executable or documented route, output hash, tolerance, and falsifier.\\
Source audit & Predecessor/source tradition named with bounded influence on OC claims.\\
Strict salvage & Workflow and reproducibility evidence without automatic Core promotion.\\
Owner gate & Explicit no-send terminal state and no outbound action.\\
\end{longtable}

\clearpage
\section*{Technical Appendix C: Decision Boundary}
The release is technically strong when its evidence topology is inspectable. It does not need to solve every future theorem to be a good 1.3.2 baseline. It does need to prevent every unsupported strong phrase from becoming public authority. That is the purpose of the claim ledger, support map, closure ledger, and release gates working together.

\end{document}
